% Round 16. Activated after independent review by global_height_audit and model_obstruction.
% Reviewed mathematical source SHA256: 0959e289975e1662723a62f800be3464e983bca8f73fa0bbb56cbe8ca0aabdfb.
% This update changes only these status comments.
\section{Principal-cycle push and a height bridge from strong G itself}
\label{r16ph:section}

This section separates two assertions which must not be confused.
The actual obstruction-zero criterion remains only a sufficient
criterion for strong G.  Nevertheless the principal transformation
law for a height can be obtained without proving the converse and
without a pullback theorem for an arbitrary supplied Green form.
The key is to push a scalar principal Green function, and to evaluate
that particular construction by the already proved point-slice residue
of a supported-complex curved pair.

\begin{lemma}[The direct immersion construction allows a normal carrier]
\label{r16ph:normal-immersion}
Let $i:M\hookrightarrow P$ be a closed immersion with $M$ normal
integral projective and $P$ smooth integral projective over $K$, of
codimension $c>0$.  Suppose an ordinary actual strong Green form is
given for a specified codimension-$r$ cycle $D$ on $M$, with
$1\le r\le\dim M$, against all compact actual tests on $M$.
Then the construction of Theorem~\ref{r16gp:push} gives an ordinary
actual strong Green form for $i_*D$ on its exact complement in $P$.
It may use the particular supported-complex curved pair for $[M]$.
\end{lemma}

\begin{proof}
Use the frozen arbitrary algebraic immersion theorem
\texttt{round14/foundations\_audit/arbitrary\_algebraic\_immersion\_germs.tex},
\texttt{r14l:local-lift}, \texttt{r14l:germ-iso} and
\texttt{r14l:global-extension}.  Its hypotheses are reduced pure
algebraic varieties; neither smoothness nor normality of the closed
source is required.  Thus a degreewise extension of the supplied form
exists on $U=P^{\rm an}\setminus|i_*D|^{\rm an}$.  The extension's
$D$ is zero on an ordinary neighborhood of $M\cap U$.  The locally
finite ordinary smooth cutoff argument of
Lemma~\ref{r16gp:open-extension} uses only paracompactness of this
smooth ambient open, so is unchanged.

The frozen supported-complex curved pair exists for the arbitrary
integral cycle $[M]$ on the smooth target $P$.  In particular its
existence does not assert that $M$ is a regular immersion.  Localize
this pair exactly as in GP16.1.  The product
$\Omega_N\widetilde\gamma_D$ is actual on all of $U$.
For an ambient complementary test $\eta$, the actual product
$\widetilde\gamma_D D\eta$ is compact on $U$, its $D$ is zero
on the support of the localized pair, and its restriction to $M$
is $\gamma_D D(i^*\eta)$.  GP16.2 therefore gives precisely
GP16.5.  All those statements make sense on the pure normal $M$;
actual-current injectivity and scalar Stokes on normal spaces are
the frozen interfaces.  No normal bundle of $M$ or smooth retraction
from $P$ to all of $M$ is used.
\end{proof}

\begin{definition}[Principal datum with a neutral omitted fiber]
\label{r16ph:datum}
Let $X/K$ be smooth projective geometrically integral of dimension
$n$, and put $p+q=n+1$, $1\le p,q\le n$.
A principal datum consists of a normal integral projective
$q$-fold $M$, a morphism $f:M\to X$, and a morphism
$u:M\to\mathbf P^1$ not identically $0$ or $\infty$.
Put $D=\operatorname{div}(u)$, $Z=f_*D$, and $S=f(|D|)$.
Let $W$ be a rational codimension-$q$ cycle on $X$.  We require
\[
 |W|\cap S=\varnothing,\qquad
 f^{-1}|W|\setminus u^{-1}(-1)\ \hbox{is finite}.
 \tag{PH16.1}\label{r16ph:position}
\]
The second condition permits arbitrary intersection in the omitted
fiber.  It is the condition supplied by a good higher-Chow
one-chain with parameter space $\mathbf P^1\setminus\{-1\}$.
All rational coefficients are cleared before forming a field norm.
\end{definition}

\begin{proposition}[An actual principal push with its exact norm evaluation]
\label{r16ph:principal-push}
For a datum in Definition~\ref{r16ph:datum}, there is an ordinary
actual form $g_Z$ of degree $(p-1,p-1)$ on
$X^{\rm an}\setminus S^{\rm an}$, with zero-curvature Green
identity for the specified cycle $Z=f_*D$ against every compact
actual test, such that
\[
                \int_W g_Z=-\log|\mu(M,u;W)|_K.
 \tag{PH16.2}\label{r16ph:norm-evaluation}
\]
For integral data, $\mu$ is the product of the field norms of $u$
at the refined zero-dimensional intersections with $W$ outside
$u=-1$, with their refined integer coefficients.  The refinement
is the graph intersection with the smooth target $X$; it is not
replaced by the length of an underived intersection on a possibly
non-Cohen--Macaulay carrier.  For rational data the displayed
logarithm is divided by the common integral scaling.
If $S=|Z|$, this is an exact-support strong Green form for $Z$.
No removability across $S\setminus|Z|$ is asserted or needed.
\end{proposition}

\begin{proof}
Write $W$ as its finite sum of prime components.  Make every
neighborhood and cutoff choice below simultaneously for that finite
union, constructing one form $g_Z$ with the one prescribed residue
$Z$.  The displayed evaluation may then be read on each prime and
summed with its coefficient; the Green forms themselves are not
summed once for each prime of $W$.  Choose a projective
embedding $e:M\hookrightarrow\mathbf P^a$ with
$c=n+a-q>0$, and let $i=(f,e):M\hookrightarrow P=X\times\mathbf P^a$.
The scalar $g=-\log|u|$ is an ordinary actual strong Green function
for $D$ on the normal $M$, by scalar all-actual Poincare--Lelong.
Let $U=P^{\rm an}\setminus|i_*D|^{\rm an}$.
Extend $g$ degreewise to an actual scalar $\widetilde g$ on $U$.
Its $D$ is zero near $i(M\setminus|D|)$ by the preceding lemma.
Choose the particular supported-complex pair $(\Gamma_M,\Omega_M)$
for the cycle $i(M)$, and a localized pair
\[
 \Gamma_N=\lambda\Gamma_M,\qquad
 \Omega_N=\Omega_M+D((\lambda-1)\Gamma_M)
 \tag{PH16.3}\label{r16ph:localized}
\]
on $U$, with support in a neighborhood where $D\widetilde g=0$.
The neighborhood and cutoff will be made smaller below.
Then $\Omega_N\widetilde g$ is the strong Green form for $i_*D$
in Lemma~\ref{r16ph:normal-immersion}.  Set
\[
 g_Z=I_{P|_{X\setminus S}/(X\setminus S)}
                  (\Omega_N\widetilde g).
 \tag{PH16.4}\label{r16ph:bundle-output}
\]
Its degree is $c-a=p-1$.  The proof of DP16.4 uses only the
just proved normal-carrier immersion result, projective-bundle
actual adjunction, and proper cycle integration.  It therefore
proves its zero-curvature identity with residue $f_*D$ and polar
envelope $S$.  In particular no proper actual operator on $f$
has been presumed.

Here is the evaluation, including the omitted fiber.
Let $W^\nu\to W$ be the normalization and put
$T=W^\nu\times\mathbf P^a$, with its proper map $b:T\to P$.
This is normal integral of dimension $c$.  Its native top integral
is that on $W\times\mathbf P^a$, since normalization has generic
degree one.  The cycle restriction formula PG16.4, applied on
the complement of $S$, gives
\[
          \int_W g_Z=\int_T b^*(\Omega_N\widetilde g).
 \tag{PH16.5}\label{r16ph:top-evaluation}
\]
This use of PG16.4 is legal on an ordinary neighborhood of the
compact cycle $W\subset X\setminus S$: its smooth-locus proof
and proper fiber support are local on that neighborhood.

Write $B=i(u^{-1}(-1))$.  The function $g$ is identically zero
on an ordinary neighborhood of $u^{-1}(-1)$, because
$|u+1|<1$ implies $|u|=1$.  The arbitrary immersion germ theorem
therefore makes $\widetilde g$ zero on an ordinary ambient
neighborhood $O_B$ of the compact $B$.  This neighborhood is
contained in $U$.  On
$T^\circ=T\setminus b^{-1}B$, the inverse image of $i(M)$
is finite by \eqref{r16ph:position}.  Denote its refined point
cycle by
\[
                         I_b([M])=\sum_x m_x[x].
 \tag{PH16.6}\label{r16ph:points}
\]
There are finitely many corresponding ambient closed points
$t=i(m)$ outside $B$.  At each such point the restriction of
$\widetilde g$ to $i(M)$ equals $-\log|u(m)|$ on a neighborhood:
$u$ is a unit there and the absolute value of a unit is locally
constant near a rigid point.  The immersion germ theorem makes
the same constant equality hold on an ambient ordinary neighborhood
of $t$.  Choose these neighborhoods disjoint and disjoint from $B$.
Choose ordinary smooth ambient functions $\chi_t$ supported in
them and equal to one on smaller neighborhoods of $t$.
Shrink $O_B$ inside the zero neighborhood of $\widetilde g$ so
that it is disjoint from every $\operatorname{Supp}\chi_t$.
This is possible because the finite union of those closed
supports is disjoint from the compact $B$.

Shrink $N$ so that its intersection with the compact set $b(T)$
lies either in $O_B$, where $\widetilde g=0$, or in those smaller
constant neighborhoods, where the corresponding $\chi_t=1$.
This is possible because the remaining closed subset of $b(T)$
is disjoint from $i(M\setminus|D|)$; ordinary paracompactness
gives the required cutoff $\lambda$.  Consequently the integral
in \eqref{r16ph:top-evaluation} is
\[
       \sum_t -\log|u(m_t)|
                       \int_{T^\circ}b^*\Omega_N\,b^*\chi_t.
 \tag{PH16.7}\label{r16ph:localized-evaluation}
\]
All displayed tests have compact support in $T^\circ$.
Use H16.4 for the particular pair $(\Gamma_M,\Omega_M)$
and the normal testing map $T^\circ\to P$.  The correction
$(\lambda-1)\Gamma_M$ is actual on $U$, so its ordinary pullback
and Stokes give the same point-slice formula for the localized
pair \eqref{r16ph:localized}.
The derivative $D(b^*\chi_t)$ is zero on the support of
$b^*\Gamma_N$ where it can contribute, by the chosen support
conditions.  Thus
\[
           \int_{T^\circ}b^*\Omega_N\,b^*\chi_t
                     =\sum_{b(x)=t}m_x[\kappa(x):K].
 \tag{PH16.8}\label{r16ph:point-mass}
\]
There is no residue calculation on the possibly larger set over
$B$; its contribution has already vanished as an actual germ.
Equations \eqref{r16ph:localized-evaluation}--\eqref{r16ph:point-mass}
and the identity
$\log|N_{\kappa(x)/K}a|=[\kappa(x):K]\log|a|$
prove \eqref{r16ph:norm-evaluation}.

Finally, proper bivariant projection identifies the push of
\eqref{r16ph:points} with the refined intersection of $M\to X$
and $W\to X$ outside $u=-1$.  To see the compatibility, use the
graph $M\hookrightarrow X\times\mathbf P^a$, flat pullback
of $W$ to the projective factor, and then proper push by the
projection.  Refined graph intersection commutes with that flat
pullback and proper push; normalizing $W$ pushes its fundamental
cycle to $[W]$.  These are the same bivariant identities used
in H16.3 (Stacks 0FBP, 0FBR and 0FC0).  The point groups on the
finite supports are free, so this keeps their coefficients.
Transitivity of field norms gives exactly the stated $\mu$,
independently of the embedding and normalization.
\end{proof}

\begin{corollary}[Principal transformation from an arbitrary strong Green form]
\label{r16ph:principal-law}
Assume in addition that $W$ has any ordinary actual strong Green
form $\gamma_W$ on its exact complement.  For a sum of principal
data as above, with $Z=\sum f_*\operatorname{div}(u)$, one has
\[
                   \int_Z\gamma_W
                       =-\log|\mu(\xi;W)|_K.
 \tag{PH16.9}\label{r16ph:principal-law-eq}
\]
The same formula gives the difference between local pairings of
two disjoint representatives $Z_0,Z_1$ connected by a good
higher-Chow one-chain $\xi$, provided their strong Green forms
and that of $W$ are supplied.  No P property is required for
the arbitrary form $\gamma_W$.
\end{corollary}

\begin{proof}
The polar envelope $S$ in Proposition~\ref{r16ph:principal-push}
is disjoint from $|W|$, and has dimension at most $n-p$.
The frozen disjoint-polar actual Stokes identity therefore applies
to $g_Z$ and $\gamma_W$, with the specified residue $Z$ even when
its support is smaller than $S$.  It gives
$\int_Z\gamma_W=\int_W g_Z$.
Use \eqref{r16ph:norm-evaluation} and add the terms of the chain.
For the difference statement, the same Stokes identity applied to
the two supplied forms gives
$\lambda(Z_1,W)-\lambda(Z_0,W)=\int_{Z_1-Z_0}\gamma_W$.
It uses their separate polar supports, both disjoint from $W$;
it does not require extending their difference across canceled
poles.  The normalization of the projective closure of a good
one-chain is a normal projective $q$-fold with its morphism to
$\mathbf P^1$.  Proper intersection on the parameter open
gives the finite set in \eqref{r16ph:position}, and the individual
faces avoid $W$.  The omitted parameter is precisely $-1$.
Thus the proposition applies without resolving the carrier.
The norm is the usual higher-Chow biextension transition norm,
up to factors of absolute value one from the omitted value and
sign convention, as in the proof of Theorem~\ref{r16h:global-height}.
\end{proof}

\begin{theorem}[A general strong-G hypothesis suffices for the full height construction]
\label{r16ph:general-height}
Let $X/F$ be smooth projective geometrically integral over a
number field, with $\dim X=n$ and $p+q=n+1$.
Assume that, for every finite place $v$, strong G holds over
$F_v$ for every specified geometrically homologically trivial
rational cycle of codimension $p$ or $q$ on $X_{F_v}$.
Then there is a canonical real bilinear height on the full groups
\[
 \operatorname{CH}^p(X)_{\mathbf Q,\mathrm{geom\,hom}}
 \times\operatorname{CH}^q(X)_{\mathbf Q,\mathrm{geom\,hom}}
                         \longrightarrow\mathbf R,
 \qquad
 h([Z],[W])=\sum_v\frac{[F_v:\mathbf Q_v]}{[F:\mathbf Q]}
                         \int_W\gamma_{Z,v},
 \tag{PH16.10}\label{r16ph:general-height-formula}
\]
with the fixed classical principal normalization at the infinite
places.  It is independent of the disjoint representatives and
all Green choices, symmetric with the complementary degrees
interchanged, and compatible with finite number-field extension
on pulled-back classes (and with the full height over the extension
whenever its corresponding strong-G hypothesis is also supplied).
It is the height of the Bloch biextension equipped with the
local norms prescribed by these integrals.

The theorem is a conditional implication from strong G.  It does
not prove its hypothesis for a general $X$, nor infer
$G(Z)\Rightarrow\mathfrak o_X([Z])=0$.
\end{theorem}

\begin{proof}
Move the two classes to disjoint representatives over $F$.
The strong-G hypothesis supplies the finite-place forms; the
usual compact K\"ahler construction supplies the archimedean ones.
The frozen disjoint-polar Stokes identity proves local Green-choice
independence and symmetry.  At all but finitely many finite places,
the pair has disjoint closures on a smooth projective model and
the frozen R11 good-model supported-complex construction has zero
curvature and vanishes outside the prescribed tube.  Its local
pairing is zero.  Choice independence therefore makes the displayed
sum finite for every permitted choice of Green forms.

A good higher-Chow chain over $F$, moved properly with respect to
the other cycle, connects two rationally equivalent disjoint
representatives.  Corollary~\ref{r16ph:principal-law} gives the
finite-place change as $-\log|\mu|_v$ for the single global
transition element after clearing all rational denominators.
At infinity the classical Green-function principal law gives the
same expression.  The weighted product formula makes its sum zero.
Symmetry handles the second variable; a third common moving
representative handles simultaneous changes.  Additivity on cycles
and this descent give rational bilinearity.  The algebraic
higher-Chow cocycle identities are precisely the Bloch-biextension
transition identities, so exponentiation gives its local norms.
These algebraic inputs and the full moving statement are those
checked in Gorchinskiy, Section 2 and Section 4.1, in the earlier
height proof.

For finite field extensions, use native finite actual trace to
see that the pullback of a strong Green form is strong on all
actual tests: in the pairing, trace the test first and apply
finite-flat cycle restriction with the full generic lengths.
The actual sheaves for a finite extension viewed over the original
field have the same genuine unit presenters, and Mihatsch
Lemma 3.16 scales their native measures by the full degree.
The extension formula for local absolute values and
$\sum_{w\mid v}[E_w:F_v]=[E:F]$ give equality of the normalized
global sums.  This repeats the native calculation in the earlier
height theorem with arbitrary strong forms; it does not use
obstruction vanishing.  The hypothesis is only used to supply
forms on all representatives, so no unjustified exact-support
addition or removability is hidden in the Chow descent.
\end{proof}

\paragraph{What this bridge removes, and what remains.}
The proof of the principal law uses the P theorem for a specific
curved pair of an auxiliary cycle, which has been proved in this
round.  It never applies P to an arbitrary strong Green form.
Consequently the unresolved converse from G to the actual
obstruction kernel is not an extra premise for the general
height implication.  Supplied regular-model comparison is a
separate theorem of this report.  Existence of strong forms for
general homological cycles remains the substantive unresolved
local hypothesis.
